\newcommand{\thmNormInducesMetric}{%
If $(X,\norm{\cdot})$ is a normed space, then
$d(x,y):=\norm{x-y}$
defines a metric on $X$.\label{thm:norm-induces-metric}%
}
\newcommand{\thmClosureSequential}{%
Let $(X,d)$ be a metric space and let $A\subset X$ be nonempty. Then:
\begin{enumerate}
\item $x\in\ol A$ if and only if there exists a sequence $(x_n)\subset A$ such that $x_n\to x$;
\item $A$ is closed if and only if every convergent sequence in $A$ has its limit in $A$.
\end{enumerate}\label{thm:closure-sequential}%
}
\newcommand{\thmSubspaceCompleteIffClosed}{%
Let $Y$ be a subspace of a Banach space $X$. Then $Y$ is complete with the inherited norm if and only if $Y$ is closed in $X$.\label{thm:subspace-complete-iff-closed}%
}
\newcommand{\thmContinuityPreimages}{%
Let $f:X\to Y$ be a map between topological spaces. Then $f$ is continuous if and only if $f^{-1}(U)$ is open in $X$ for every open set $U\subset Y$.\label{thm:continuity-preimages}%
}
\newcommand{\thmOperatorNorm}{%
Let $X,Y$ be normed spaces. For a linear map $L:X\to Y$, define
$\norm{L}_{\mathcal L(X,Y)}:=\sup_{\norm{x}_X\le 1}\norm{Lx}_Y
=\sup_{\norm{x}_X=1}\norm{Lx}_Y.$
Then this defines a norm on the space $\mathcal L(X,Y)$ of bounded linear maps from $X$ to $Y$.\label{thm:operator-norm}%
}
\newcommand{\thmBoundedIffContinuous}{%
For a linear map $L:X\to Y$ between normed spaces, the following are equivalent:
\begin{enumerate}
\item $L$ is bounded;
\item $L$ is continuous at $0$;
\item $L$ is continuous at one point of $X$;
\item $L$ is continuous on $X$.
\end{enumerate}\label{thm:bounded-iff-continuous}%
}
\newcommand{\thmDualIsBanach}{%
If $Y$ is Banach, then $\bigl(\mathcal L(X,Y),\norm{\cdot}_{\mathcal L(X,Y)}\bigr)$ is Banach. In particular, for every normed space $X$, its dual $X^*=\mathcal L(X,\K)$ is Banach.\label{thm:dual-is-banach}%
}
\newcommand{\thmFiniteDimRange}{%
If $\dim X=n<\infty$ and $L:X\to Y$ is linear, then $\dim \operatorname{Ran}L\le n$.\label{thm:finite-dim-range}%
}
\newcommand{\thmFiniteDimBounded}{%
If $\dim X<\infty$, then every linear map $L:X\to Y$ is bounded.\label{thm:finite-dim-bounded}%
}
\newcommand{\thmFiniteDimProperties}{%
Let $X$ be a normed space with $\dim X=N<\infty$. Then:
\begin{enumerate}
\item $X$ is linearly homeomorphic to $\K^N$;
\item $X$ is Banach;
\item any two norms on $X$ are equivalent.
\end{enumerate}\label{thm:finite-dim-properties}%
}
\newcommand{\thmHamelBasis}{%
Every vector space has a Hamel basis.\label{thm:hamel-basis}%
}
\newcommand{\thmProperSubspace}{%
If $X$ is finite-dimensional and $Y\subsetneq X$ is a proper subspace, then $\dim Y<\dim X$.\label{thm:proper-subspace}%
}
\newcommand{\thmCompactEquivSequential}{%
For a subset $K$ of a metric space, the following are equivalent:
\begin{enumerate}
\item $K$ is compact;
\item every sequence in $K$ has a convergent subsequence with limit in $K$;
\item $K$ is complete and totally bounded.
\end{enumerate}\label{thm:compact-equiv-sequential}%
}
\newcommand{\thmUnitBallCompactIffFiniteDim}{%
A normed space $X$ is finite-dimensional if and only if its closed unit ball
$\Ball=\{x\in X:\norm{x}\le 1\}$
is compact.\label{thm:unit-ball-compact-iff-finite-dim}%
}
\newcommand{\thmFrechetMetric}{%
Let $(p_k)_{k\in\N}$ be a separating sequence of seminorms on a vector space $X$. Then
$d(x,y):=\sum_{k=1}^{\infty}2^{-k}\frac{p_k(x-y)}{1+p_k(x-y)}$
defines a metric on $X$. If this metric is complete, then $X$ is a Frechet space.\label{thm:frechet-metric}%
}
\newcommand{\thmHahnBanachReal}{%
Let $X$ be a real vector space, let $p:X\to\R$ be sublinear, let $V\subset X$ be a subspace, and let $f:V\to\R$ be linear with $f(v)\le p(v)$ for all $v\in V$. Then there exists a linear $F:X\to\R$ such that $F|_V=f$ and $F(x)\le p(x)$ for all $x\in X$.\label{thm:hahn-banach-real}%
}
\newcommand{\thmHahnBanachComplex}{%
Let $X$ be a complex vector space, let $p:X\to\R$ satisfy $p(x+y)\le p(x)+p(y)$ and $p(\lambda x)=\abs{\lambda}p(x)$, let $V\subset X$ be a subspace, and let $f:V\to\C$ be linear with $\abs{f(v)}\le p(v)$ for all $v\in V$. Then there exists a linear $F:X\to\C$ such that $F|_V=f$ and $\abs{F(x)}\le p(x)$ for all $x\in X$.\label{thm:hahn-banach-complex}%
}
\newcommand{\thmHahnBanachNormed}{%
Let $V$ be a subspace of a normed space $X$ and let $f\in V^*$. Then there exists $F\in X^*$ such that $F|_V=f$ and $\norm{F}_{X^*}=\norm{f}_{V^*}$.\label{thm:hahn-banach-normed}%
}
\newcommand{\thmHahnBanachSeparating}{%
Let $Y\subset X$ be a closed subspace of a normed space and let $x_0\notin Y$. Then there exists $\phi\in X^*$ such that
$\phi|_Y=0, \qquad \phi(x_0)=\operatorname{dist}(x_0,Y), \qquad \norm{\phi}=1.$\label{thm:hahn-banach-separating}%
}
\newcommand{\thmDualSeparates}{%
For every nonzero $x\in X$ there exists $\phi\in X^*$ such that $\phi(x)\ne 0$. Consequently weak convergence is well-defined.\label{thm:dual-separates}%
}
\newcommand{\thmNormViaDual}{%
For every $x$ in a normed space $X$,
$\norm{x}=\sup\{\abs{\phi(x)}: \phi\in X^*,\ \norm{\phi}\le 1\}.$\label{thm:norm-via-dual}%
}
\newcommand{\thmDualSeparableImpliesSeparable}{%
If $X$ is Banach and $X^*$ is separable, then $X$ is separable.\label{thm:dual-separable-implies-separable}%
}
\newcommand{\thmBanachAlaoglu}{%
Let $X$ be a normed space. The closed unit ball of $X^*$ is compact in the weak-* topology. In particular, if $X$ is separable, every bounded sequence in $X^*$ has a weak-* convergent subsequence.\label{thm:banach-alaoglu}%
}
\newcommand{\thmReflexiveWeaklyCompact}{%
If $X$ is reflexive, then the closed unit ball of $X$ is weakly compact. Equivalently, every bounded sequence in $X$ has a weakly convergent subsequence.\label{thm:reflexive-weakly-compact}%
}
\newcommand{\thmCantorIntersection}{%
Let $(X,d)$ be complete and let $F_1\supset F_2\supset \cdots$ be nonempty closed sets with $\operatorname{diam}(F_n)\to 0$. Then $\bigcap_{n=1}^{\infty}F_n$ consists of exactly one point.\label{thm:cantor-intersection}%
}
\newcommand{\thmBaireCategory}{%
A complete metric space cannot be a countable union of nowhere dense sets. Equivalently, a countable intersection of dense open sets is dense.\label{thm:baire-category}%
}
\newcommand{\thmBaireVariant}{%
If a complete metric space is the union of countably many closed sets, then at least one of those closed sets has nonempty interior.\label{thm:baire-variant}%
}
\newcommand{\thmNoCountableHamelBasis}{%
An infinite-dimensional Banach space cannot have a countable Hamel basis.\label{thm:no-countable-hamel-basis}%
}
\newcommand{\thmBanachSteinhaus}{%
Let $X$ be Banach, $Y$ normed, and $\mathcal F\subset\mathcal L(X,Y)$. If for every $x\in X$,
$\sup_{L\in\mathcal F}\norm{Lx}_Y<\infty,$
then
$\sup_{L\in\mathcal F}\norm{L}_{\mathcal L(X,Y)}<\infty.$\label{thm:banach-steinhaus}%
}
\newcommand{\thmLimitOperatorBounded}{%
Let $X$ be Banach, $Y$ normed, and let $(L_n)\subset \mathcal L(X,Y)$. If for every $x\in X$ the limit
$Lx:=\lim_{n\to\infty}L_nx$
exists, then $L\in\mathcal L(X,Y)$.\label{thm:limit-operator-bounded}%
}
\newcommand{\thmOpenMapping}{%
Let $X,Y$ be Banach spaces and let $L\in\mathcal L(X,Y)$ be surjective. Then $L$ maps open sets in $X$ onto open sets in $Y$.\label{thm:open-mapping}%
}
\newcommand{\thmBoundedInverse}{%
Let $X,Y$ be Banach spaces and let $L\in\mathcal L(X,Y)$ be bijective. Then $L^{-1}\in\mathcal L(Y,X)$.\label{thm:bounded-inverse}%
}
\newcommand{\thmComparableNorms}{%
Let $\norm{\cdot}_1$ and $\norm{\cdot}_2$ be norms on the same vector space $X$. Assume both $(X,\norm{\cdot}_1)$ and $(X,\norm{\cdot}_2)$ are Banach and that there exists $C>0$ with
$\norm{x}_2\le C\norm{x}_1 \qquad\text{for all }x\in X.$
Then the two norms are equivalent.\label{thm:comparable-norms}%
}
\newcommand{\thmClosedGraph}{%
Let $X,Y$ be Banach spaces and let $L:X\to Y$ be linear. If
$\operatorname{Graph}(L)=\{(x,Lx):x\in X\}$
is closed in $X\times Y$, then $L\in\mathcal L(X,Y)$.\label{thm:closed-graph}%
}
\newcommand{\thmClosedOperatorFacts}{%
Let $D(L)\subset X$ and let $L:D(L)\to Y$ be linear.
\begin{enumerate}
\item $L$ is closed if and only if its graph is closed in $X\times Y$.
\item If $D(L)$ is closed in $X$ and $L:D(L)\to Y$ is bounded with respect to the inherited norm on $D(L)$, then $L$ is closed.
\item The assumption $D(L)=X$ is essential in the closed graph theorem.
\end{enumerate}\label{thm:closed-operator-facts}%
}
\newcommand{\thmAdjoint}{%
Let $L\in\mathcal L(X,Y)$ and define $L^*:Y^*\to X^*$ by $L^*\psi=\psi\circ L$. Then $L^*$ is linear and bounded, $\norm{L^*}=\norm{L}$, and $(S\circ L)^*=L^*\circ S^*$ whenever the composition is defined.\label{thm:adjoint}%
}
\newcommand{\thmCompactLimit}{%
If $K_n\in\mathcal L(X,Y)$ are compact and $K_n\to K$ in operator norm, then $K$ is compact.\label{thm:compact-limit}%
}
\newcommand{\thmCompactWeakToStrong}{%
If $K\in\mathcal L(X,Y)$ is compact and $x_n\rightharpoonup x$ weakly in $X$, then $Kx_n\to Kx$ strongly in $Y$.\label{thm:compact-weak-to-strong}%
}
\newcommand{\thmSchauder}{%
A bounded linear operator $K:X\to Y$ is compact if and only if its adjoint $K^*:Y^*\to X^*$ is compact.\label{thm:schauder}%
}
\newcommand{\thmIntegralOperator}{%
Let $k\in C([a,b]\times[a,b])$ and define
$(Kf)(x):=\int_a^b k(x,y)f(y)\,dy \qquad (f\in C([a,b])).$
Then $K:C([a,b])\to C([a,b])$ is compact.\label{thm:integral-operator}%
}
\newcommand{\thmInnerProductNorm}{%
If $(\cdot,\cdot)_H$ is an inner product on a vector space $H$, then $\norm{x}_H:=\sqrt{(x,x)_H}$ defines a norm.\label{thm:inner-product-norm}%
}
\newcommand{\thmCauchySchwarz}{%
For all $x,y\in H$,
$\abs{(x,y)_H}\le \norm{x}_H\norm{y}_H.$\label{thm:cauchy-schwarz}%
}
\newcommand{\thmParallelogram}{%
For all $x,y\in H$,
$\norm{x+y}_H^2+\norm{x-y}_H^2=2\norm{x}_H^2+2\norm{y}_H^2.$\label{thm:parallelogram}%
}
\newcommand{\thmBestApprox}{%
Let $e_1,\dots,e_n$ be orthonormal in $H$. Then for every $x\in H$ and every scalars $c_1,\dots,c_n$,
$\norm{x-\sum_{k=1}^n c_ke_k}_H^2
=
\norm{x-\sum_{k=1}^n (x,e_k)_He_k}_H^2
+\sum_{k=1}^n \abs{c_k-(x,e_k)_H}^2.$
Hence $\sum_{k=1}^n (x,e_k)_He_k$ is the best approximation to $x$ from $\operatorname{span}\{e_1,\dots,e_n\}$.\label{thm:best-approx}%
}
\newcommand{\thmBesselParseval}{%
If $(e_k)_{k\ge 1}$ is orthonormal in $H$, then
$\sum_{k=1}^{\infty}\abs{(x,e_k)_H}^2\le \norm{x}_H^2$
for every $x\in H$.
Equality for every $x$ holds if and only if $(e_k)$ is complete; then
$x=\sum_{k=1}^{\infty}(x,e_k)_He_k,
\qquad
\norm{x}_H^2=\sum_{k=1}^{\infty}\abs{(x,e_k)_H}^2.$\label{thm:bessel-parseval}%
}
\newcommand{\thmEllTwoIsometry}{%
Let $(e_k)_{k\ge 1}$ be orthonormal in $H$. Then the map
$\ell^2\ni (a_k)\longmapsto \sum_{k=1}^{\infty} a_ke_k$
is an isometry onto $\ol{\operatorname{span}}\{e_k:k\in\N\}$.\label{thm:ell-two-isometry}%
}
\newcommand{\thmOrthonormalCompleteEquiv}{%
For an orthonormal system $(e_k)$ in $H$, the following are equivalent:
\begin{enumerate}
\item $(e_k)$ is complete;
\item if $(x,e_k)=0$ for all $k$, then $x=0$;
\item the span of $(e_k)$ is dense in $H$;
\item Parseval's identity holds for every $x\in H$.
\end{enumerate}\label{thm:orthonormal-complete-equiv}%
}
\newcommand{\thmSeparableHilbertONB}{%
Every separable Hilbert space has a complete orthonormal system.\label{thm:separable-hilbert-onb}%
}
\newcommand{\thmSeparableHilbertIsomorphicEllTwo}{%
Every separable Hilbert space is isometrically isomorphic to $\ell^2$.\label{thm:separable-hilbert-isomorphic-ell-two}%
}
\newcommand{\thmProjection}{%
Let $H$ be a Hilbert space and let $C\subset H$ be nonempty, closed, and convex. Then for every $x\in H$ there exists a unique $y\in C$ such that
$\norm{x-y}_H=\inf_{z\in C}\norm{x-z}_H.$
If $C=M$ is a closed subspace, then this minimizer is characterized by $x-y\in M^{\perp}$.\label{thm:projection}%
}
\newcommand{\thmDirectSum}{%
If $M$ is a closed subspace of a Hilbert space $H$, then
$H=M\oplus M^{\perp}.$\label{thm:direct-sum}%
}
\newcommand{\thmLeastSquares}{%
Let $A\in\R^{m\times n}$ and $b\in\R^m$. A vector $x\in\R^n$ minimizes $\norm{Ax-b}_2$ if and only if it solves the normal equations
$A^TAx=A^Tb.$\label{thm:least-squares}%
}
\newcommand{\thmRieszRepresentation}{%
Let $H$ be a Hilbert space. For every $\phi\in H^*$ there exists a unique $a\in H$ such that
$\phi(x)=(x,a)_H \qquad\text{for all }x\in H,$
and $\norm{\phi}_{H^*}=\norm{a}_H$.\label{thm:riesz-representation}%
}
\newcommand{\thmLaxMilgramSymmetric}{%
Let $H$ be a real Hilbert space, and let $B:H\times H\to\R$ be bounded, symmetric, and coercive:
$\abs{B(u,v)}\le C\norm{u}\norm{v}, \qquad B(u,u)\ge \beta\norm{u}^2.$
Then for every $f\in H^*$ there exists a unique $u\in H$ with
$B(u,v)=f(v) \qquad \forall v\in H.$
Moreover $u$ is the unique minimizer of
$J(v)=\frac12 B(v,v)-f(v).$\label{thm:lax-milgram-symmetric}%
}
\newcommand{\thmCoerciveBijective}{%
Let $H$ be a Hilbert space and let $L\in\mathcal L(H)$. Assume there exists $\beta>0$ such that
$\operatorname{Re}(Lu,u)_H\ge \beta\norm{u}_H^2 \qquad\text{for all }u\in H.$
Then $L$ is bijective and
$\norm{L^{-1}}\le \frac{1}{\beta}.$\label{thm:coercive-bijective}%
}
\newcommand{\thmLaxMilgramGeneral}{%
Let $H$ be a Hilbert space over $\K$, let $B:H\times H\to\K$ be sesquilinear, bounded, and coercive:
$\abs{B(u,v)}\le C\norm{u}\norm{v}, \qquad \operatorname{Re}B(u,u)\ge \beta\norm{u}^2.$
Then for every $f\in H^*$ there exists a unique $u\in H$ such that
$B(u,v)=f(v) \qquad \forall v\in H,$
and $\norm{u}\le \beta^{-1}\norm{f}$.\label{thm:lax-milgram-general}%
}
\newcommand{\thmHilbertAdjoint}{%
Let $H$ be a Hilbert space and let $T\in\mathcal L(H)$. Then there exists a unique $T^*\in\mathcal L(H)$ such that
$(Tx,y)_H=(x,T^*y)_H \qquad\text{for all }x,y\in H.$
Moreover $\norm{T^*}=\norm{T}$ and $(ST)^*=T^*S^*$.\label{thm:hilbert-adjoint}%
}
\newcommand{\thmCompactWeakToStrongHilbert}{%
If $K\in\mathcal L(H)$ is compact and $x_n\rightharpoonup x$ weakly in $H$, then $Kx_n\to Kx$ strongly in $H$.\label{thm:compact-weak-to-strong-hilbert}%
}
\newcommand{\thmFredholmAlternative}{%
Let $H$ be a Hilbert space and let $K\in\mathcal L(H)$ be compact. Then:
\begin{enumerate}
\item $\ker(I-K)$ and $\ker(I-K^*)$ are finite-dimensional;
\item $\operatorname{Ran}(I-K)=\ker(I-K^*)^{\perp}$;
\item $\dim\ker(I-K)=\dim\ker(I-K^*)$;
\item either $\ker(I-K)=\{0\}$ and $(I-K)^{-1}\in\mathcal L(H)$, or $\ker(I-K)$ is nontrivial and solutions of $(I-K)u=f$ (when they exist) form the affine space $u_0+\ker(I-K)$.
\end{enumerate}\label{thm:fredholm-alternative}%
}
\newcommand{\thmNeumannSeries}{%
Let $T\in\mathcal L(X)$ with $\norm{T}<1$. Then $I-T$ is invertible and
$(I-T)^{-1}=\sum_{n=0}^{\infty}T^n$
with convergence in operator norm.\label{thm:neumann-series}%
}
\newcommand{\thmResolventAnalytic}{%
Let $X$ be a complex Banach space and $T\in\mathcal L(X)$. The resolvent set
$\rho(T)=\{\lambda\in\C: \lambda I-T \text{ is invertible}\}$
is open, and the resolvent map $\lambda\mapsto R(\lambda,T)=(\lambda I-T)^{-1}$ is analytic on $\rho(T)$.\label{thm:resolvent-analytic}%
}
\newcommand{\thmSpectralRadius}{%
Let $X\neq\{0\}$ be a complex Banach space and let $T\in\mathcal L(X)$. Then $\sigma(T)$ is nonempty and compact, and
$r(T):=\sup_{\lambda\in\sigma(T)}\abs{\lambda}
=\lim_{m\to\infty}\norm{T^m}^{1/m}.$
If $\Gamma$ is a positively oriented contour surrounding $\sigma(T)$, then
$T^j=\frac{1}{2\pi i}\int_{\Gamma}\lambda^jR(\lambda,T)\,d\lambda.$\label{thm:spectral-radius}%
}
\newcommand{\thmCompactSpectrum}{%
If $K\in\mathcal L(H)$ is compact on a Hilbert space $H$, then every nonzero point of $\sigma(K)$ is an eigenvalue. The set $\sigma(K)\setminus\{0\}$ is at most countable, each nonzero eigenvalue has finite multiplicity, and its only possible accumulation point is $0$.\label{thm:compact-spectrum}%
}
\newcommand{\thmHellingerToeplitz}{%
Every everywhere-defined symmetric linear operator on a Hilbert space is bounded.\label{thm:hellinger-toeplitz}%
}
\newcommand{\thmSelfadjointSpectrumReal}{%
If $T\in\mathcal L(H)$ is selfadjoint, then $\sigma(T)\subset\R$. More precisely,
$\sigma(T)\subset [m,M], \quad m:=\inf_{\norm{x}=1}(Tx,x), \quad M:=\sup_{\norm{x}=1}(Tx,x).$\label{thm:selfadjoint-spectrum-real}%
}
\newcommand{\thmHilbertSchmidt}{%
Let $H$ be a \textbf{separable} Hilbert space and let $T\in\mathcal L(H)$ be compact and selfadjoint. Then there exists an orthonormal basis $(e_n)$ of $H$ consisting of eigenvectors of $T$, with real eigenvalues $\lambda_n\to 0$, such that
$Tx=\sum_{n=1}^{\infty}\lambda_n(x,e_n)e_n \qquad\text{for all }x\in H.$\label{thm:hilbert-schmidt}%
}
